% Chapter 5: early results, with a pgfplots figure and a small table.
\chapter{Preliminary results}\label{ch:prelim}

A first version of the model was fitted to the 214 municipal sensors
operating in Aldmere between 2021 and 2025. Leave-one-out validation on
summer afternoons gives an RMSE of 0.82\,K and a mean 90\,\% interval
coverage of 0.88, close to nominal. Figure~\ref{fig:curve} compares
placement rules in a retrospective simulation: starting from 20 random
sensors, each rule adds sensors from the remaining pool and the error is
measured on all withheld sensors, weighted by exposure.

\begin{figure}[ht]
\centering
\begin{tikzpicture}
\begin{axis}[
  width=.92\linewidth, height=58mm,
  xmin=20, xmax=200, ymin=0.6, ymax=2.0,
  xlabel={Sensors placed}, ylabel={Weighted RMSE (K)},
  axis lines=left, axis line style={muted, line width=.4pt},
  tick style={muted}, tick label style={font=\sffamily\footnotesize},
  label style={font=\sffamily\small, text=muted},
  ymajorgrids, grid style={muted!20, line width=.3pt},
  legend style={draw=none, font=\sffamily\footnotesize, at={(1,1)},
    anchor=north east, fill=none}, legend cell align=left,
  every axis plot/.append style={line width=1pt, mark=none, smooth}]
\addplot[muted!70, dashed] coordinates {(20,1.92) (40,1.66) (60,1.47)
  (80,1.33) (100,1.22) (120,1.13) (140,1.06) (160,1.00) (180,0.96) (200,0.93)};
\addlegendentry{Random}
\addplot[accentdark!55] coordinates {(20,1.92) (40,1.52) (60,1.32)
  (80,1.19) (100,1.09) (120,1.02) (140,0.97) (160,0.93) (180,0.90) (200,0.88)};
\addlegendentry{Variance only}
\addplot[accent, line width=1.6pt] coordinates {(20,1.92) (40,1.36) (60,1.12)
  (80,0.99) (100,0.91) (120,0.85) (140,0.81) (160,0.78) (180,0.76) (200,0.74)};
\addlegendentry{Exposure weighted}
\draw[accent, densely dotted, line width=.5pt] (axis cs:20,1.13) -- (axis cs:200,1.13);
\fill[accent] (axis cs:70,1.13) circle (1.3pt);
\fill[muted] (axis cs:120,1.13) circle (1.3pt);
\node[font=\sffamily\scriptsize, text=accent, anchor=south] at (axis cs:70,1.14) {70};
\node[font=\sffamily\scriptsize, text=muted, anchor=north] at (axis cs:120,1.12) {120};
\end{axis}
\end{tikzpicture}
\caption{Retrospective placement in Aldmere, mean of 50 random starts.
The dotted line marks the error random placement reaches with 120
sensors; the weighted rule reaches it with about 70.}
\label{fig:curve}
\end{figure}

Table~\ref{tab:prelim} summarises the same experiment at 100 sensors.
The weighted rule gives up little unweighted accuracy, which supports
the trade-off assumed in H1, though the simulation draws only from
sites that already host a sensor and so understates the choice
available in a real deployment.

\begin{table}[ht]
\centering\sffamily\small
\caption{Placement rules at 100 sensors (Aldmere, summer afternoons).}
\label{tab:prelim}
\begin{tabular}{@{}lrrr@{}}
\toprule
\thead{Rule} & \thead{Weighted RMSE} & \thead{Unweighted RMSE}
  & \thead{Coverage}\\
\midrule
Random            & 1.22 K & 0.97 K & 0.86 \\
Variance only     & 1.09 K & 0.89 K & 0.88 \\
\textcolor{accent}{Exposure weighted} & \textcolor{accent}{0.91 K}
  & 0.92 K & 0.89 \\
\bottomrule
\end{tabular}
\end{table}
